\documentclass

\usepackage{asmmath}
\usepackage{amssymb}
\usepackage{geometry}

\geometry{margin=1in}

\begin{document}

\begin{flushright}
CSE 21 Fa26
Homework 2
Submitted: 10/12/2026
\end{flushright}

\begin{center}
David Washington(200484943), Judah Boyce(200453537), Osinachi Agaranna(200455712), Nick Young(200455733)
\end{center}

\begin{enumerate}

\item
\begin{enumerate}

\item How many 7 characte license plates start with a letter or end with a numeral?

Let
\[
L = \text{Plates that start with a letter}
\]
\[
N = \text{plates that end with a numeral}
\]

We want to find $|L \cup N|$.

Using the inclusion-exclusion principle, 

\[
|L \cup N| = |L| + |N| - |L \cap N|
\]

There are 26 letter, 10 numerals, and 36 possible character in all.

For plates that start with a letter, the first character has 26 possibilites and the remaining 6 have 36 possibilites.

\[
|L| = 26 \cdo 36^6
\]

For plates that end with a numeral, the frist 6 character have 36 possibilites and the last has 10 possibilites.

\[
|N| = 36^6 \cdot 10
\]

For plates that start with a letter and end with a numeral, there are 26 choices for the first character, 10 for the last, and 36 for each of the 5 middle characters.

\[
|L \cap N| = 26 \cdot 36^5 \cdot 10
\]

Therefore,
\[
\boxed{
|L \cup N|
=
26 \cdot 36^6
+ 
10 \cdot 36^6
-
26 \cdot 10 \cdot 36^5
}
\]

\item How many 7 character license plates use at least one A, at least one B, and at least one numeral?

Let
\[
A = \text{Plates without an A}
\]

\[
B = \text {Plates without a B}
\]

\[
N = \text{Plates without a numeral}
\]

We want a plate containing at least one A, one B, and one numeral.

We use the complement and the inclusion-exclusion principle. There are $36^7$ possible plates. The individual sets have sizes 

\[
|A| = 35^7, \qquad |B| = 35^7
\]

\[
|N| = 26^7
\]

Their pairwise intersections are 

\[
|A \cap B| = 34^7
\]

\[
|A \cap N| = 25^7
\]

\[
|B \cap N| = 25^7
\]

The interesection of all three sets is 

\[
|A \cap B \cap N| = 24^7
\]

Using inclusion-exclsuion

\begin{align*}
|A \cup B \cup N|
&= |A| + |B| +|N|\\
&\quad - |A \cap B|\\
&\quad - |A \cap N|\\
&\quad - |B \cap N|\\
&\quad + |A \cap B \cap N|
\end{align*}

We subtract this from the total possibilites:

\[
\boxed{
36^7 -
\left(
2(35^7) + 26^7
- 34^7 - 2(25^7) + 24^7
\right)
}
\]

\item How many 7 character license plates contain only numerals that add up to 10?

We need to the find the number of solutions to 

\[
x_1 + x_2 + \cdots + x_7 = 10
\]

where 

\[
0 \leq x_i \leq 0
\]

We can use the stars and bars formula to count the solutions

\[
\binom{n+k-1}{k-1}
\]

Here, $n=10$ and $k=7$.

\[
\binom{10+7-1}{7-1} 
=
\binom{16}{6}
\]

However, stars and bars allows individual values greater than 9. 

Since the toal sum is 10, the only invalid cases are there where one digit equals 10
and the remaining 6 equal 0.

For example,

\[
(10,0,0,0,0,0,0)
\]

\[
(0,10,0,0,0,0,0)
\]

\[
\vdots
\]

There are exactly 7 invalid cases, one for each possible position of 10. We subtract these cases:

\[
\boxed{
\binom{16}{6}-7
}
\]

\item How many 7 character license plates contain only letters in alphabetical order, with repitions allowed?

We use combination with repition:

\[
\binom {n+k-1}{k}
\]

where 
\[
n = 26 \text { letters}
\]

\[
k = 7 \text { characters}}
\]

Because the letters must appear in alphabetical order, each selection of 7 letters corresponds to
exactly one valid license plate.

Substituting

\[
\binom{26+7-1}{7}
=
\binom{32}{7}
\]

Therefore, 

\[
\boxed{
\binom{32}{7}
}
\]

\end{enumerate}

\item
\begin{enumerate}

\item How many different ways cam the museum curator arrange all 9 prints among the 12 rooms?

We have 9 identical prints and 12 different rooms.

We use stars and bars because we are distrbuting identical objects among distinct groups.

Let $x_i$ represent the number of prints placed in room $i$. 

We need

\[
x_1 + x_2 + \cdots + x_{12} = 9
\]

where $x_i \geq 0$.

Using stars and bars, 

\[
\binom{n+k-1}{k-1}
\]

where $n=9$ and $k=12$.

\[
\binom{9+12-1}{12-1}
=
\binom{20}{11}
\]

Therefore,

\[
\boxed{
\binom{20}{11}
}
\]


\item How many ways can the museum curator arrange all 9 prints among the 12 rooms so that each
print is in a different room?

Since all 9 prints are identical and must be placed in different rooms, we only need to choose which 9 of the 
12 rooms recieve a print, regardless of order.

Therefore,

\[
\boxed{
\binom{12}{9}
}
\]


\item How many ways can the museum curator arrange all 9 prints among the 12 rooms so that at least 3 rooms are non-empty?

We can  use complementary countin.

We begin with all possibile distributions and subtract the ones where only 1 or 2 rooms are non-empty.

From part (a), the total number of distributions is:

\[
\binom{20}{11}
\]

First, consider exactly 1 non-empty room.

We chhose 1 of the 12 rooms to contain all 9 prints

\[
\binom{12}{1} =12
\]

Next consider exactly 2 non-empty rooms.

We first chhose 2 rooms from the 12

\[
\binom{12}{2}
\]

We then distribute the 9 identical prints between those two 

Let $x_1$ and $x_2$ represent the number of prints in each room.

\[
x_1 + x_2 = 9 
\]

where $x_1, x_2 \geq 1$.

using stars and bars:

\[
\binom{9-1}{2-1}
= 
\binom{8}{1} 
\]

Thus, the number of distribution with exactly 2 nonempty rooms is 

\[
\binom{12}{2}]binom{8}{1}
\]

Subtracting both invalid cases from the total gives

\[
\boxed{
\binom{20}{11}
-
!2
-
\binom{12}{2}\binom{8}{1}
}
\]

% -----------2d goes here--------
/item How many ways can the museum curator arrange all 9 prints among the 12 rooms such that at
least one room has exactly 3 prints?

We use the inclusion-exclusion principle.

Let $A_i$ be the set of distributions where room $i$ contains exactly 3 prints.

We want to find

\[
\left|\bigcup_{i=1}^{12} A_i\right|
\]

First, consider one room containing exactly 3 prints. 

We choose 1 of the 12 rooms.

The remaining 6 prints are distributed among the 11 other rooms.

Using stars and bars,

\[
\binom{6+11-1}{11-1}
=
\binom{16}{10}
\]

Thus, the first term is

\[
\binom{12}{1}\binom{16}{10}
\]

Next, consider two rooms that each contain exactly 3 prints.

We choose 2 of the 12 rooms.

The remaining 3 prints are distributed among the other 10 rooms.

Using stars and bars,

\[
\binom{3+10-1}{10-1}
=
\binom{12}{9}
\]

Thus, the second term is 

\[
\binom{12}{2}\binom{12}{9}
\]

Finally, consider three rooms that contain exactly 3 prints.

Since all 9 prints have been used the remaining rooms are empty

We choose 3 of the 12 rooms

\[
\binom{12}{3}
\]

There can't be four rooms with exactly 3 prints because that would require 12 prints.

Using inclusion-exclusion, we add the single room cases, subtract double room, and add back the triple room cases:

\[
\boxed{
\binom{12}{1}\binom{16}{10}
-
\binom{12}{2}\binom{12}{9}
+
\binom{12}{3}
}
\]

\item How many ways can the museum curator arrange all 9 prints in 3 of the 12 rooms such that each room has exactly 3 prints?

We need to choose exactly 3 rooms from the available 12.

Since the 9 prints are identical, once we select the 3 rooms, there is only one way to place exactly 3 prints in each room.

Therefore, we use combination to choose the 3 rooms:

\[
\boxed{
\binom{12}{3}
}
\]

\end{enumerate}





\item 
We want to find the number of solutions to 

\[ a_1+a_2+a_3+a_4+a_5+a_6=30.
\]

\begin{enumerate}
  \item
  We are given the restrictions
  \[a_1\geq1,\ a_2\geq2,\ a_3\geq3,\ a_4\geq4,\ a_5\geq5,\ a_6\geq6.
  \]

  First, assign each variable its minimu required value.
  \[
  1+2+3+4+5+6=21.
  \]

  This leave $30-21=9$ remaining. Let $x_i=a_i-i$ for $1\leq i\leq6$.
  Then
  \[
  x_1+x_2+x_3+x_4+x_5+x_6=9,
  \]
  where each $x_i\geq0$.

  Using 9 stars and 5 bars the number of solutions is:

  \[
  \boxed{\binom{9+6-1}{6-1}=\binom{14}{5}}.
  \]

  \item
  We want at least one variable to be less than or equal to 3.

  We use the complement method by subtracting the solution where every variable is at least 4 from the total number of solution

  Using Stars and Bars, the total number of solution without any restriction is:

  \[ \binom{30+6-1}{6-1}=\binom{35}{5}.
  \]

  For the complement, each variable must satisfy $a_i\geq4$. Assigning 4 to each variable uses:

  \[ 
  6*4=24,
  \]
  leaving $30-24=6$.

  Let $x_i=a_i-4$. Then
  \[
  x_1+x_2+x_3+x_4+x_5+x_6=6,
  \]
  where each $x_i\geq0$

  By Stars and Bars, the complement solutins is:

  \[
  \binom{6+6-1}{6-1}=\binom{11}{5}.
  \]

  Subtracting these from the total gives the valid solutions:
  \[
  \boxed{\binom{35}{5}-\binom{11}{5}}.
  \]

  \item
  We are given the restrictions 
  \[
  0\leq a_i\leq7
  \quad \text{for all}1\leq i\leq6.
  \]

  Use Inclusion-Exclusion to subtract the solutions where at least one variable is greater than or equal to 8.

  No restrictions, using Stars and Bars:

  \[
  \binom{35}{5}
  \]
  solutions.

  Let $A_i$ be the set of solutions where $a_i\geq8$.

  If one variable is at least 8, we assign 8 to that variable, leaving $30-8=22$.
  There are $\binom{6}{1}$ ways to select the variable:

  \[
  \binom{6}{1}\binom{22+5}{5}
  = binom{6}{1}\binom{27}{5}
  \]

  Now if two variables are at least 8 we are left with $30-16=14$.
  There are $\binom{6}{2}$ ways:
  \[
  \binom{6}{2}\binom{14+5}{5}
  =\binom{6}{2}\binom{19}{5}.
  \]

  Moving on to if three variable are at least 8 we are left with $30-24=6$.
  There are $\binom{6}{3}$ ways:
   \[
  \binom{6}{3}\binom{6+5}{5}
  =\binom{6}{2}\binom{11}{5}.
  \]

  This is the max case because if four variables are at least 8 we get  $4*8=32>30$.

  Therefore the number of valid solutions is given by:

  \[
  \boxed{
  \binom{35}{5}
  -\binom{6}{1}\binom{27}{5}
  +\binom{6}{2}\binom{19}{5}
  -\binom{6}{3}\binom{11}{5}
  }.
  \]
\end{enumerate}
